%% Este template foi desenvolvido com base no AbnTex2 para a criação de projetos de AIPA destinados ao IPHAN. Foi adicionado um pacote que traduz diversos comandos do LaTeX comumente utilizados para o português (ver exemplos na Introdução). No entanto, você pode continuar utilizando os comandos do LaTeX normalmente, caso prefira.
%%
%% O mantenedor atual deste template é Geraldo Pereira de M. Jr.
%% Contato: geraldo.pmj@gmail.com
%%
\documentclass[
	12pt,				   % tamanho da fonte
	openright,             %
	oneside,			   % 
	a4paper,			   % tamanho do papel.
	sumario = tradicional, %
	english,			   % idioma adicional para hifenização
	french,				   % idioma adicional para hifenização
	spanish,		       % idioma adicional para hifenização
	brazil			       % o último idioma é o principal do documento
	]{abntex2}
\input{Documento/config/packages}%Alterar apenas se necessário
% ---
% Informações para adicionar automaticamente no projeto:
% ---
%Sobre a empresa de arqueologia
\newcommand{\empresaarqueologia}{Arcehocy Atrarmolchatio}% Nome da empresa de arqueologia
\newcommand{\localempresaarqueologia}{Belo Horizonte}% Local da empresa de arqueologia
\newcommand{\coordenadorgeral}{Fulano de Alguem }% Nome do coordenador geral
\newcommand{\coordenadordecampo}{Ciclano de Tals }% Nome do coordenador de campo
%
% Sobre o processo e empreendimento
\newcommand{\processoiphan}{00000.000000/ano-mês}% Número do processo Iphan
\newcommand{\portaria}{Nº x, de dia de mês de ano}% Portaria
\newcommand{\apoioinstitucional}{Reserva Técnica}% Nome da reserva técnica
\newcommand{\empreendimento}{Nome do Empreendimento}% Nome do empreendimento
\newcommand{\contratante}{Empresa Tal LTDA}% Nome da empresa contratante
\newcommand{\localidade}{Rio de Janeiro - RJ}% Local de execução do empreendimento
\newcommand{\enquadramentodoprojeto}{Nível III}% Nível de acordo com IN 01/2015
%
\input{Documento/config/capaEcontracapa}%Alterar apenas se necessário
% ---
% compila o indice - Não mexer
% ---
\makeindex
% ---

% ----
% Início do documento
% ----
\begin{document}
\selectlanguage{brazil}
% Retira espaço extra obsoleto entre as frases.
\frenchspacing 
%----------------------------------------------------
% ELEMENTOS PRÉ-TEXTUAIS
%----------------------------------------------------
\pretextual %- Não mexer
%---
% Capa
%---
\imprimircapa %- Não mexer
%---
% Folha de rosto
% (o * indica que haverá a ficha bibliográfica)
%---

\imprimirfolhaderosto* %- Não mexer

%---
%inserir o sumario 
%---
\input{capitulos/sum.tex} %- Não mexer

%---
%LISTAS
%---
\input{capitulos/Listas.tex} %- Não mexer

%-----------------------------------------------------
% ELEMENTOS TEXTUAIS
%-----------------------------------------------------

\textual %- Não mexer

%-----------------------------------------------------
%-----------------------------------------------------
\input{Documento/capitulos/0-introducao}
%-----------------------------------------------------
\input{Documento/capitulos/1-LocalizacaoEcaracterizacao}
%-----------------------------------------------------
\input{Documento/capitulos/2-Contexto arqueologico e etnohistorico}
%-----------------------------------------------------
\input{Documento/capitulos/3-Metodologia}
%-----------------------------------------------------
\input{Documento/capitulos/5-Cronograma de execução}
%-----------------------------------------------------
\input{Documento/capitulos/7-Equipe}
%-----------------------------------------------------

%------------------------------------------------------
% ELEMENTOS PÓS-TEXTUAIS 
%------------------------------------------------------

\postextual %- Não mexer
\pagestyle{plain} %- Não mexer

%----------------------------------------------------
% Referências bibliográficas
%----------------------------------------------------

\printbibliography[title={Referências Bibliográficas}] %- Não mexer

%-----------------------------------------------------
% Apêndices 
%-----------------------------------------------------

\begin{apendicesenv}
\input{capitulos/apend.tex}
\end{apendicesenv}

%-----------------------------------------------------
% Anexos
%-----------------------------------------------------

%\begin{anexosenv}
%\input{capitulos/anex.tex}
%\end{anexosenv}

%------------------------------------------------------
%-----------------------------------------------------
\end{document}
